% ===== 原文 pp.1–2（PDF pp.2–3）=====

% [原文 p.1]
\section{引言}
\label{sec:intro}

本文是 2016 年 8 月 26--28 日在东京大学数理科学研究科（Graduate School of Mathematical Sciences,
the University of Tokyo）所作讲座的讲义的一部分。完整版本将很快面世。

\section{二维 $(2,2)$ 超对称量子场论}
\label{sec:2d}

我们首先介绍二维 $(2,2)$ 超对称量子场论（QFT）的基础，着重于其数学层面。

\subsection{二维 $(2,2)$ 超对称}
\label{sec:22susy}

在 Minkowski 时空上表述时，一个二维 $(2,2)$ 超对称 QFT 具有下述对称性算子：时间平移 $H$
（Hamiltonian）、空间平移 $P$（动量）、Lorentz 变换 $M$、超荷 $Q_+,\overline Q_+,Q_-,\overline Q_-$，
以及可能存在的矢量 R-荷 $F_V$ 和/或轴 R-荷 $F_A$。它们作用在 $\Z_2$ 分次的态空间 $\cH$ 上。
超荷 $Q_\pm$ 与 $\overline Q_\pm$ 是奇的、且互为伴随，其余算子则是偶的、自伴的。
关于 Lorentz 群，$H$ 与 $P$ 构成一个矢量，$i[M,H\pm P]=\mp 2(H\pm P)$；超荷是旋量，
$i[M,Q_\pm]=\mp Q_\pm$，$i[M,\overline Q_\pm]=\mp\overline Q_\pm$；R-荷是标量，
$[M,F_V]=[M,F_A]=0$。超荷满足\footnote{对 (2.2) 可能存在由中心项给出的修正，本讲义不予考虑。}
\begin{equation}
  \{Q_\pm,\overline Q_\pm\}=H\pm P, \tag{2.1}
\end{equation}
\begin{equation}
  \text{其余所有反对易子}=0. \tag{2.2}
\end{equation}

R-荷是超荷的相位旋转：
\begin{align}
  [F_V,Q_\pm]&=-Q_\pm, & [F_V,\overline Q_\pm]&=\overline Q_\pm, \tag{2.3}\\
  [F_A,Q_\pm]&=\mp Q_\pm, & [F_A,\overline Q_\pm]&=\pm\overline Q_\pm. \tag{2.4}
\end{align}

\subsection{$A$ 与 $B$}
\label{sec:AandB}

令
\begin{equation}
  Q_A:=\overline Q_++Q_-,\qquad Q_B:=\overline Q_++\overline Q_-. \tag{2.5}
\end{equation}
% [Tutor Note: 这里 Q_A、Q_B 的取法并非任意：(2.1) 给出 {Q_±,\overline Q_±}=H±P，(2.2) 给出其余
% 反对易子为零，于是 Q_A²={Q_-,Q_-}+2{\overline Q_+,Q_-}+{\overline Q_+,\overline Q_+}=0，
% Q_B²=2{Q_+,\overline Q_-}+2{\overline Q_+,\overline Q_-}=0 都自动成立；再配合 (2.3)(2.4)，
% [F_A,Q_A]=Q_A、[F_V,Q_B]=Q_B 也成立，因此 (Q_A,F_A) 与 (Q_B,F_V) 各自给出一个以 Q²=0 为
% 微分的复形。注意 Q_A 与 Q_B 都含 \overline Q_+，这是本讲义采用的约定。]

% [原文 p.2]
则 $(Q,F)=(Q_A,F_A)$ 或 $(Q_B,F_V)$ 满足
\begin{equation}
  Q^2=0,\qquad [F,Q]=Q. \tag{2.6}
\end{equation}
这意味着态空间构成一个以 $Q$ 为微分、以 $F$ 为分次的复形。\footnote{本小节中我们的论断假定了
$F_A$ 或 $F_V$ 存在，但这并非必需：我们也可以用 $\Z_2$ 分次代替。}%
局部算子的空间也同理。特别地，局部算子的上同调类构成一个环，称为\emph{手征环}：
对应 $(Q_A,F_A)$ 的记为 $\cR_A$，对应 $(Q_B,F_V)$ 的记为 $\cR_B$。它是一个分次交换代数。

当表述在半 Minkowski 空间上时——例如空间坐标被限制为 $x\le 0$、时间 $t$ 无界——必须在 $x=0$
处为场指定一个边界条件。本质上存在两种保持最多超荷数的边界条件：
\[
\begin{array}{ll}
A\text{-型：} & Q_A\ \text{与}\ Q_A^\dagger\ \text{守恒};\\[2pt]
B\text{-型：} & Q_B\ \text{与}\ Q_B^\dagger\ \text{守恒}.
\end{array}
\]
这类边界条件分别称为 \emph{A-膜} 与 \emph{B-膜}。配对 $(Q,F)=(Q_A,F_A)$（相应地 $(Q_B,F_V)$）
作用在插入于具有 A-型（相应地 B-型）边界条件的边界上的局域算子，并满足与 (2.6) 相同的关系。
其上的上同调类构成一个代数，一般是非交换的。两个同型的边界条件可以置于边界上，并在二者之间
插入一个局域算子。配对 $(Q,F)$ 同样作用在这样插入的局域算子上并满足 (2.6)，\footnote{这些关系
[1] 可能存在一个潜在的异常。} 于是我们可以考虑其上同调类。这样我们便得到范畴：A-膜的范畴记为
$\cC_A$，B-膜的范畴记为 $\cC_B$。对象是边界条件；边界条件之间的态射是插入二者之间的局域算子的
$Q$-上同调类，复合则由算子的乘积表示。

\subsection{重整化群流}
\label{sec:RG}

在一般 QFT 中，可观测量的行为极其依赖于能标，或者反过来说，依赖于距离标度。例如，插入在相距
$r$ 的两点处的两个算子的关联函数，一般是 $r$ 的一个复杂函数。小 $r$（即短距离，或等价地，高能、
紫外 (UV)）处的行为，一般与大 $r$（即长距离，或等价地，低能、红外 (IR)）处的行为大不相同。
若增大长度标度（即降低能标），
